\documentclass[10pt]{article}
\usepackage[margin=20mm]{geometry}
\usepackage{amsmath,amssymb,mathtools,fontspec,unicode-math,xurl,hyperref}
\IfFontExistsTF{DejaVu Serif}{\setmainfont{DejaVu Serif}}{\setmainfont{Latin Modern Roman}}
\setmathfont{latinmodern-math.otf}
\setlength{\parindent}{0pt}\setlength{\parskip}{6pt}\sloppy
\title{Initial scientific review of Intrinsic mirror symmetry}
\author{}
\date{}
\begin{document}
\maketitle
Abstract: This claim-focused review concerns the acquired author manuscript arXiv:1909.07649v3, not an independently compared journal version. A false assertion about geometric injectivity of saturation occurs in the proof of Lemma 8.4, on the dependency path to associativity. An explicit family of monoids disproves that assertion. The offending points are outside the actual closed monomial fibre: a direct argument restores the particular uniqueness statement needed there. The supported judgment is Major for the demonstrated core proof-step defect, with no demonstrated counterexample to the lemma or the main theorems. The checked scope comprises 90 of 175 text pages and 68 visually inspected original pages. This is not a whole-paper correctness certification.\par
\section*{Identity, version, and scope}
Paper M25-0039 is Mark Gross and Bernd Siebert, Intrinsic mirror symmetry, DOI https://doi.org/10.1090/jams/1066. The reviewed PDF is https://arxiv.org/pdf/1909.07649v3, submitted 30 September 2025, with internal date 1 October 2025, 175 PDF pages. Page references below use the PDF’s one-based page numbers, which agree with its printed numbering. The locally acquired released source is the corresponding v3 source tree. In particular, PDF Section 8 is in the file `section7.tex`.\par
The source manifest identifies this as an author manuscript. The journal version was not acquired or compared; the source-preparation record reports an AMS rate limit. The official arXiv abstract record was also consulted for version identity. Its version description is not proof of equality with journal pagination or text. This report makes no claim that the identified issue persists in the journal version, and does not claim an exhaustive later-version or erratum search.\par
Text actually read: pp. 1–24, 35–38, 51–61, 69–83, 93–100, 105–106, and 132–157 (90 pages). Original images actually inspected: pp. 1, 7, 10–16, 19–23, 35–37, 51–54, 57–60, 69, 73–75, 81–83, 93–100, 105–106, and 132–157 (68 pages). The exact machine-readable lists are in `coverage.json` and `evidence/visual\_inspection.json`. Extraction of all pages and rendering 90 pages are recorded separately and are not counted as reading or visual inspection.\par
\section*{Definitions, quantifiers, and checked structure}
The setup is over an algebraically closed field of characteristic zero. Basic Setup 1.1 requires a projective logarithmically smooth map to a point or to the specified non-complete logarithmic curve, with the stated Zariski and global-generation hypotheses. Basic Setup 1.6 includes the curve-class monoid, finite-type stable-map conditions, and the co-Artinian monomial ideal. These assumptions are substantive inputs to finiteness and properness.\par
Theorem 1.9 allows the logarithmic first Chern class to be nef or anti-nef. Theorem 1.15 instead uses the log Calabi–Yau setup, the nonnegative boundary coefficients for the negative logarithmic first Chern class, and the subcomplex \(B\) of good rays. It does not replace \(B(\mathbb Z)\) with all of \(\Sigma(X)(\mathbb Z)\). The structure constants use a logarithmic point constraint and a negative output contact order; they are defined to vanish unless \(A\cdot c_1(\Theta_{X/k})=0\).\par
The principal checked dependency is: \[\text{Lemma 8.4(1)}
\longrightarrow \text{Theorem 8.12}
\longrightarrow \text{Theorem 4.28}
\longrightarrow \text{Theorem 4.4}
\longrightarrow \text{Theorems 1.9 and 1.15}.\] Theorem 8.12, Step 1 (p. 152), explicitly invokes Lemma 8.4(1). Theorem 4.4 on p. 83 combines the comparisons and independence of modulus. Other dependencies in that chain have only the extent of verification described below.\par
\section*{Problem 1: saturation is not globally injective on geometric points}
\subsection*{Location and exact logical issue}
Equation reference: eq:normalization. Lemma 8.4, proof Step 3, p. 142, introduces \(P=Q_\lambda^{\mathrm{fine}}=Q\oplus_R^{\mathrm{fine}}R_\lambda\subset Q^\mathrm{gp}\). It then reasons from torsion-freeness of \(Q^\mathrm{gp}\), and hence of the saturation \(P^\mathrm{sat}\), to geometric injectivity of \[\operatorname{Spec}k[P^\mathrm{sat}]\longrightarrow\operatorname{Spec}k[P].
\tag{1}\] This implication is false. The same statement is present in released `section7.tex`, lines 451–454, so it is not a PDF extraction error. Step 4 uses Step 3 to establish pure-dimensionality in Lemma 8.4(1).\par
\subsection*{Complete countercheck of the asserted affine-chart implication}
Equation reference: eq:explicit. Take \(k=\mathbb C\), fix integers \(\mu\geq1\) and \(N\geq1\), and set \[Q=\mathbb Ne_1+\mathbb Ne_2,\quad
 \delta=e_2,\quad \ell_q=2e_1+\mu e_2,\quad \ell=\ell_q,\quad
 \lambda=\mu+2N.\] Thus \(Q\) is a sharp rank-two fine saturated monoid, \(Q^\mathrm{gp}=\mathbb Z^2\) is torsion-free, and \(Q\) is rationally generated by \(\delta\) and \(\ell_q-\mu\delta=2e_1\), as in the nearby local model of Lemma 8.3. Let \[R=\mathbb N\delta+\mathbb N\ell,\qquad
 R_\lambda=\mathbb N\delta+\mathbb N(\ell-\lambda\delta)\subset R^\mathrm{gp}.\] These are fine saturated free monoids in their own common group: \(R^\mathrm{gp}=R_\lambda^\mathrm{gp}\). Use the homomorphism \(\theta:R\to Q\) with the displayed images. It has the form (8.2) with \(a=1,b=0\). The integral image of the pushout, hence the fine pushout, is \[P=\mathbb Ne_1+\mathbb Ne_2+\mathbb N(2e_1-2Ne_2)\subset\mathbb Z^2.\] Its group is \(\mathbb Z^2\). It is finitely generated, integral and sharp: the functional \((a,b)\mapsto(N+1)a+b\) is positive on each nonzero displayed generator. Its saturation is \[P^\mathrm{sat}=\mathbb Ne_2+\mathbb N(e_1-Ne_2).\] Indeed these two vectors are a lattice basis, of determinant \(-1\), and generate the rational cone of \(P\). Writing \(x=\chi^{e_1}\), \(y=\chi^{e_2}\), \(z=\chi^{2e_1-2Ne_2}\), and \(t=\chi^{e_1-Ne_2}\), one obtains \[k[P]=k[x,y,z]/(x^2-y^{2N}z),\qquad
 k[P^\mathrm{sat}]=k[y,t],\qquad
 x\mapsto y^Nt,\quad y\mapsto y,\quad z\mapsto t^2.
\tag{2}\] For completeness, the relation is the primitive lattice relation on the three generators; the rank-two semigroup algebra has this prime height-one relation ideal. The element \(t\) is integral (\(t^2=z\)) and lies in the fraction field (\(t=x/y^N\)); the polynomial ring \(k[y,t]\) is normal. This also verifies the normalization description directly.\par
The distinct geometric points \((y,t)=(0,1)\) and \((0,-1)\) both map to \((x,y,z)=(0,0,1)\). Thus (1) is not one-to-one, even though all the stated local algebraic conditions above hold. Taking \(N\) arbitrarily large makes \(\lambda-\mu\) arbitrarily large. A requirement that \(\lambda\) be sufficiently large does not establish the false affine-chart assertion.\par
\subsection*{Counterevidence and the exact scope of the counterexample}
This is a counterexample to the unrestricted affine-chart assertion in Step 3. It is not an example of two punctured stable maps in the actual fibre of \(i\), and does not refute Lemma 8.4. The counterexample point has \(z=1\), whereas the new positive generator of \(R_\lambda\) vanishes on the closed monomial stratum used in \(B\mathbb G_m^\ddagger\). The paper’s geometric application therefore has a further restriction which excludes precisely these points.\par
The surrounding argument fixes a geometric ghost stalk, so the disputed sentence admits the intended reading that uniqueness is asserted only over its closed characteristic point. That restricted statement is valid by the argument below. The finding concerns the explicitly written global affine-chart premise; it does not establish a failure of the actual geometric-fibre conclusion.\par
This restriction is enough to supply a direct replacement argument. It must be stated explicitly instead of asserting injectivity of (1) everywhere. The distinction between the ambient chart and its closed monomial fibre is essential to both the diagnosis and the repair.\par
\subsection*{Verified replacement for the needed geometric-fibre statement}
Claim 1.  Let \(P\) be a sharp fine monoid embedded in a torsion-free group, and let \(P^\mathrm{sat}\) denote its saturation in \(P^\mathrm{gp}\). The closed monomial point of \(\operatorname{Spec}k[P]\) has exactly one geometric point above it in \(\operatorname{Spec}k[P^\mathrm{sat}]\).\par
Proof. The saturation is sharp. If \(q\in P^\mathrm{sat}\setminus\{0\}\), then \(mq\in P\setminus\{0\}\) for some positive integer \(m\). At any geometric lift, \((\chi^q)^m=\chi^{mq}=0\), so \(\chi^q=0\) in its residue field. Thus every nonzero monomial has a uniquely determined value. The map sending all such monomials to zero and the zero monomial to one exists because the monoid is sharp, so this determines precisely one geometric point. ◻\par
Apply this after choosing the local charts already used in Step 3 and trivializing the relevant torus torsors over the algebraically closed residue field. Nonzero elements of the ghost \(Q\) vanish at the chosen geometric point. Nonzero elements of \(R_\lambda\) likewise vanish on the closed monomial classifying stack. Hence every nonzero generator of the fine pushout has value zero. At a nonempty such point a positive generator cannot become a unit: it would have to map both to zero and to a unit. The characteristic monoid is therefore sharp as required. The identity of group completions used in Step 3 supplies the common lattice; no additional character choices arise from a changed group. The claim gives the unique geometric lift required for that fibre.\par
Together with the finite-surjective saturation and reduced-surjectivity argument already established in Step 3, this gives its intended bijection and supplies the input to Step 4. It proves only the needed geometric-fibre statement. It does not say the underlying nonreduced morphism is an isomorphism, nor that every multiplicity is preserved.\par
For example, the displayed family has \(K=(y,x^2)\subset k[Q]\), giving \(k[x]/(x^2)\), of length two. The corresponding saturated quotient is \(k[t]/(t^2)\), also of length two, but the induced map sends \(x=y^Nt\) to zero. Thus equality of these lengths is compatible with a nonisomorphic nonreduced morphism, exactly as the caution in footnote 8 on p. 140 allows. General multiplicity matching still uses Lemmas 8.3 and 8.10.\par
\subsection*{Impact and judgment}
The demonstrated defect is an invalid inference in a stated proof on the associativity dependency path. Under the specified core-contribution calibration its judgment is Major, proof-step defect. This judgment does not rely on an assertion that the entire mirror construction is false, that the repair is difficult, or that an offending ambient toric point is an actual punctured-map point. The closed-fibre repair above is analytically verified. Human acceptance of the finding and repair remains pending, and author adoption has not been requested. The repair is a review contribution, not an attributed author correction.\par
\section*{Other checks and exclusions}
\subsection*{Product finiteness, degree zero, and the unit}
Lemma 4.1 (p. 57) uses the balancing constraint of Remark 2.13 (pp. 35–37): the contact order has a fixed image in the space of global linear functionals. Global generation gives injectivity on each cone, and the finite cone complex gives only finitely many possibilities. The monomial quotient supplies the finite set of relevant classes. Finite units in the curve monoid do not by themselves make the coefficient quotient a field; the review does not use such an assumption.\par
For degree-zero three-point maps (Lemma 1.18, pp. 51–54), the checked argument uses the stable genus-zero three-point domain, cone balancing, quotienting logarithmic enhancements by the ghost-character action, and \(H^1(\mathbb P^1,\mathcal O(-1)^d)=0\). The asserted moduli properties still depend on the punctured-map foundations cited by the authors. Input symmetry gives commutativity.\par
In the unit argument (pp. 93–98), the evaluation correction includes the output contact term and the change in the conormal line. Forgetting a marking and pushing the zero-dimensional class to the dimension-minus-one Artin-stack target gives the required vanishing under the stated obstruction-theory compatibility. The \(A=0\) case and the convention for nonzero Chern degree are separate. The construction and properness of all preceding forgetful stacks were not independently reproved.\par
\subsection*{Tropical and virtual-dimension checks}
Proposition 4.21 (pp. 73–75) uses a two-dimensional miniversal tropical family, no terminal tail, and an unbounded length at fixed \(\delta\). The checked recurrence for consecutive vertex positions yields a nonnegative splitting direction by rescaling in a closed cone. The case \(r=0\) retains the independent parameter \(\delta\). This does not by itself exclude an internal tail.\par
For Lemma 7.17 (pp. 133–136), the review checked why the output belongs to the open stratum \(Z_r^\circ\), why \(r\) must lie in \(B(\mathbb Z)\) in the log Calabi–Yau case, and why the tail degree is nonpositive. If a positive boundary coefficient had negative intersection, the negative contact order would force the output component into that boundary divisor, contrary to the point constraint. The relative virtual dimension is \(A_1c_1(\Theta_{X/k})-\dim X\leq-\dim X\). The two-point genus-zero prestable base has dimension \(-1\); integrality and idealized logarithmic etaleness give the strict negative total virtual dimension needed for the cited vanishing theorem. The general virtual-pullback theorem in Appendix A was not independently proved in this review.\par
For Lemma 8.8 (pp. 144–151), the checked split separates non-dominating components and contact orders outside the skeleton from the remaining no-tail reduction. The finite-type restriction is essential to choose one sufficiently large \(\mu\) for finitely many cone conditions. The \(s=0\) case is explicitly accommodated by footnote 9. For Lemma 8.10 (p. 151), absence of internal tails gives coefficient one in \(\ell=\ell_q+b\delta\), the exact hypothesis needed for multiplicity matching in Lemma 8.3. Theorems 8.11–8.12 compare weighted components after tail vanishing; they do not assume all nonreduced spaces are isomorphic.\par
In Theorem 8.12 (pp. 152–157), the constraints \(\lambda\geq a+b\) and \(\lambda\geq b+\mu a\) make the chosen node vanish and the splitting factor through the modified base. The argument chooses \(\mu\) before \(\lambda\) after restricting to the finitely many relevant types (pp. 137–138 and 156). The two-sided inverse is asserted on the tail-free locus, where the splitting node is unique. The last monomorphism is a monomorphism of log stacks, justified by the group completion, not an assertion about all underlying stacks.\par
\subsection*{Examples and scope clarification}
The blow-up and line–conic examples (pp. 20–22) were checked at the level of curve intersections and possible contact orders. For the line–conic pair the equations are \(d=1-a_1,\ 2d=1-a_2\) for the first product, and \(d=2-a_1,\ 2d=2-a_2\) for the second. With nonnegative degree and coefficients these give the listed possibilities. The enumerativity of each punctured invariant is not an independent outcome of those elementary equations.\par
The relation calculation in Example 4.5 (pp. 59–60) reproduces its coefficient two algebraically. This is a check within the displayed algebraic presentation, not an independent count of maps or a proof of associativity from scratch. For \(\mathbb P^1\) with three boundary points in Example 1.27 (p. 23), the logarithmic tangent bundle has degree \(2-3=-1\), hence is anti-nef. The sign wording there is recorded as an unrated scope/typographical clarification: Theorem 1.9 already allows anti-nefness and Theorem 1.15 uses the negative Chern class. It is not a separate substantive finding.\par
\section*{Remaining work and conclusion}
Eighty-five text pages were not read. In particular, the full evaluation and virtual-smoothness foundations, all forgetful-map constructions, the intermediate adjacent-cone argument, the complete gluing proof, the remainder of Section 9, and Appendix A have not received an independent complete audit here. External punctured-map and virtual-pullback references were not fully reverified. The additional 22 rendered pages that were not displayed are not claimed as visually inspected.\par
Item | Status\par
Problem 1 | Major: false unrestricted saturation-injectivity inference.\par
Main theorem | No counterexample demonstrated; restricted-fibre repair supplied.\par
Repair | Analytic proof for the exact needed closed geometric fibre.\par
Adoption | Human acceptance pending; no author adoption requested.\par
Version | Author manuscript only; journal comparison and exhaustive errata check pending.\par
\section*{Supporting checks}
> 72e51b2658ed7df51126401b1fe3beff81da53ddb7f2163990f6e442476f6d8d.\par
The decisive released `section7.tex` SHA-256 is\par
> e6aa517c6c4f886086e18f88a19737f391de2aeb5f30c5dc9ba381a4a61f4e8b.\par
Its lines 404–469 were read as a scientific source excerpt; other TeX files were keyword-searched, not fully semantically audited. The exact excerpt is saved in `evidence/source\_excerpts.json`.\par
The reviewer’s own `evidence/independent\_checks.py` was executed with `python -X utf8`. It calculates the displayed normalization maps for \(N=1,\ldots,6\), their lattice determinants, the two distinct points with the same image, the example relation, the contact equations, and the degree \(-1\). The length-two quotient bases and closed-stratum observations are recorded analytic conclusions, not separate computer proofs. The actual process exited zero in 3.12 seconds; full results are saved in `independent\_checks.json`. These finite tests support the examples; the general counterexample family and repair are the analytic arguments given above. No author code was executed.\par
Text extraction used the existing Poppler `pdftotext -layout`; original-page rendering used `pdftoppm -f p -l p -singlefile -scale-to 1800 -png`. Exact commands, return codes, stderr and hashes are in `extraction\_log.json` and `render\_log.json`. All 90 render operations returned zero. A first-page font diagnostic was retained in the log; the inspected title/version image was readable. Only the 68 displayed pages in the separate inspection ledger count as visual QA.\par
The original compiler-absence probe preceded the authorized portable Tectonic follow-up. Its build and page inspections remain preserved as historical evidence. The current report is compiled from this exact TeX with Tectonic 0.17.0 in untrusted, cached-only mode. Exact input, dependency, log, PDF and page-image hashes, together with actual inspection of changed or reflowed pages, are recorded in the accompanying provenance. Original-review actions retain their original attribution; focused AI acceptance is recorded separately. Scientific adjudication remains pending, with no human review or author adoption claimed.\par
All scientific decisions, calculations, source inspections and original-page inspections were performed by the main review assistant. A separate agent was restricted to read-only administrative identity/path/provenance auditing. No equivalence of those labels or backend verification is claimed. No independent evidence of a backend mismatch was obtained.\par
9 M. Gross and B. Siebert, Intrinsic Mirror Symmetry, arXiv:1909.07649v3 (30 September 2025), acquired 175-page author PDF and matching released TeX source. https://arxiv.org/abs/1909.07649v3. Associated DOI: https://doi.org/10.1090/jams/1066. Local acquired-source manifest for M25-0039, `papers/0039/source\_manifest.json`; exact file identity, source version, optional-source discovery and journal-access limitations. References such as ACGS20 and Manolache12 discussed above are dependencies cited by the paper, not fully reverified external sources in this review.\par
\end{document}
